\section{The smooth-density theorem in the adaptive experiment}
\label{app:smooth-density}

This appendix formalizes \cref{thm:smooth} in the original experiment.
The policy has arbitrary measurable private randomness and may revisit
coins, interleave observations, and use an unbounded stopping time.
For an integer $s\ge1$, let $\mathcal F_{s,M}$ contain the $C^\infty$
probability densities supported in $[1/4,3/4]$ with
$\max_{k\le s}\|f^{(k)}\|_\infty\le M$.

\begin{theorem}[Verified smooth-density robustness]
\label{thm:smooth-verified}
For each $s\ge1$ there are $M_s\ge1$, $c_s>0$, and $\varepsilon_s>0$
such that, for $M\ge M_s$, $0<\varepsilon\le\varepsilon_s$, and
$0<\delta\le1/4$, the minimax expected number of actual flips satisfies
\[
 c_s\varepsilon^{-2}\log^2(1/\varepsilon)\log(1/\delta)
 \ \le\ \mathcal C_{s,M}(\varepsilon,\delta)
 \ \le\ (65536M)^2\varepsilon^{-2}\log^2(1/\varepsilon)\log(1/\delta).
\]
The upper estimator is selected before the density.
Every estimator unbiased throughout $\mathcal F_{s,M}$ has either
unbounded variance over that class or unbounded expected work.
\end{theorem}

\begin{proof}
Write $x=p-1/2$ and take the normalized even smooth bump $g$ with inner
radius $1/8$ and outer radius $3/16$. It integrates to one and satisfies
$g\ge1$ on $[-1/8,1/8]$. For $0<t\le1$, set
\[
 R(x)=\frac{x}{1+x^2},\qquad
 \eta_t(x)=\frac{t^s}{8}R(x/t),\qquad
 f_\pm(1/2+x)=g(x)(1\pm2\eta_t(x)).
\]
The perturbation is odd and bounded in absolute value by $1/8$.
Thus both densities are nonnegative and normalized.
The identity
\[
 R^{(k)}(x)=
 \operatorname{Re}\frac{(-1)^k k!}{(x+\mathrm i)^{k+1}}
\]
gives $\|R^{(k)}\|_\infty\le k!$. Consequently,
\[
 \|\eta_t^{(k)}\|_\infty
 \le \frac{k!}{8}t^{s-k}\le\frac{k!}{8},
 \qquad 0\le k\le s.
\]
Compact support bounds every fixed derivative of $g$.
Leibniz's rule then supplies one finite $M_s$ for all $t$ and both signs.

The common gap law has density $2g(d)$ on $(0,1/2]$, and its sign
probabilities are $1/2\pm\eta_t(d)$.
The formal identification tests this mixture against every bounded
measurable function of the parameter. It also identifies the threshold
probability with the original sign mean and the parameter mixture used
by the independent Bernoulli-tape simulator.

Put $h=t^{s+1}$. For every positive gap,
\[
 0\le\eta_t(d)\le\frac18,\qquad d\eta_t(d)\le\frac h8.
\]
The stopped-transcript bound therefore gives
$\operatorname{KL}(\mathbb P_+,\mathbb P_-)\le(128/225)h^2B$
whenever expected work under the plus model is at most $B$.
For $t\le1/64$, integration on $[t,1/8]$ yields
\[
 \Delta_t=\theta_+-\theta_-
 =4\int_0^\infty g(d)\eta_t(d)\,{\rm d}d
 \ge\frac h8\log(1/t).
\]
Let $L=\log(1/\varepsilon)$ and $K_s=64(s+1)$.
For sufficiently small $\varepsilon$, the choice
$h=K_s\varepsilon/L$ and $t=h^{1/(s+1)}$ has $t\le1/64$
and $\Delta_t>2\varepsilon$.
The binary testing bound is $\log(1/\delta)/4$.
Combining it with the transcript bound gives the lower result.
This construction supplies $c_s=225/(512K_s^2)$.

For the upper bound, every model in the class satisfies
$\Pr(D\le u)\le2Mu$.
The reference power law of exponent one and scale $a=(2M)^{-1}$
therefore dominates its small-gap probabilities.
At $0<z<a$ its clipped modulus is
$(z/a)(1+\log(a/z))$.
Apply the law-independent confidence estimator with
\[
 B=(65536M)^2\varepsilon^{-2}L^2\log(1/\delta),
 \qquad z=\sqrt{\log(1/\delta)/B}=\frac{\varepsilon}{65536ML}.
\]
For $\varepsilon\le e^{-16}$, its radius
$8192(z/a)(1+\log(a/z))$ is at most $\varepsilon$.

Suppose one estimator is unbiased with variance at most $V$
and expected work at most $W$ throughout the class.
The same local information bound holds in both directions.
The verified two-distribution variance inequality and $|\Delta_t|\le1$
imply
\[
 \Delta_t^2\le\frac{128}{225}h^2W(2V+1/2).
\]
Combining this with the separation bound and cancelling $h^2>0$
bounds $\log^2(1/t)$ by a constant independent of $t$.
Letting $t$ tend to zero gives a contradiction.
\end{proof}

The \href{https://github.com/davidiach/adaptive-threshold-estimation/blob/8e42a99a49dcd4af8bdaaf17a6ea1c398f34b03b/lean/ThresholdHC/SmoothDensity/Full.lean}{full Lean theorem}
checks the lower bound, the estimator family for the upper bound, and
the unbiased impossibility together. Its class definition uses equality
of the latent parameter law against all bounded measurable test functions.
The source audit checks the construction and the original-experiment
connections separately before checking this combined statement.
